%! Mutually Exclusive Events — Unit 4, Lesson 4.4 — Student Packet Cover Sheet
\documentclass[10pt]{article}
\usepackage{apstats-article}
\usepackage{apstats-boxes}
\usepackage{ltablex}
\keepXColumns

\begin{document}

% Full-bleed navy banner
\begin{tikzpicture}[remember picture, overlay]
  \fill[navy] (current page.north west)
    rectangle ([yshift=-0.9in]current page.north east);
\end{tikzpicture}
\vspace{-0.8in}

\begin{center}
  {\color{white}\textbf{\LARGE AP Statistics}}\\[4pt]
  {\color{skymid}\large\textbf{Unit 4: Probability, Random Variables, and Probability Distributions}}\\[6pt]
  {\color{white}\textbf{\large Lesson 4.4 \quad Mutually Exclusive Events}}
\end{center}

\vspace{0.22in}
\namedateperiod
\vspace{0.18in}

\begin{learningtargetbox}
\small
\begin{enumerate}[leftmargin=1.5em, itemsep=5pt]
  \item \ldots{} explain what it means for two events to be mutually exclusive (disjoint)
        and verify the condition $P(A \cap B) = 0$ in context.
  \item \ldots{} determine whether two events are mutually exclusive by checking whether
        any outcome belongs to both events simultaneously.
  \item \ldots{} justify my answer in context --- not just ``yes'' or ``no'' but
        \emph{why} two events can or cannot co-occur.
\end{enumerate}
\end{learningtargetbox}

\vspace{0.14in}

\begin{tocbox}
\small
\renewcommand{\arraystretch}{1.5}
\begin{tabularx}{\linewidth}{@{} c l X r @{}}
\rowcolor{sky}
\textbf{\#} & \textbf{Component} & \textbf{Description} & \textbf{Score} \\
\midrule
1 & Warm-Up        & Spiral review: basic probability + Venn diagram thinking   & \blank{1.2cm} \\
2 & Guided Notes   & Joint probability; mutually exclusive definition \& Venn diagrams & \blank{1.2cm} \\
3 & Group Activity & Classify event pairs; justify; compute $P(A \cap B)$        & \blank{1.2cm} \\
4 & Exit Ticket    & Streaming service two-way table                             & \blank{1.2cm} \\
5 & Homework       & Blood type \& health outcome scenarios                      & \blank{1.2cm} \\
\midrule
  & \multicolumn{2}{r}{\textbf{Total}} & \blank{1.2cm} \\
\end{tabularx}
\end{tocbox}

\vspace{0.14in}

\begin{remindbox}
\small
The \textbf{core idea of Lesson 4.4}: two events are mutually exclusive when
they \emph{cannot both happen at the same time} --- their intersection is empty.

\vspace{6pt}
\begin{tabularx}{\linewidth}{@{} >{\centering\arraybackslash\columncolor{goldbg}}p{0.06\linewidth}
                                    X @{}}
\textbf{\large!} &
\textbf{Mutually exclusive $\Rightarrow$ $P(A \cap B) = 0$.}\enspace
If two events have no outcomes in common, they are disjoint and their joint
probability is zero. \\[6pt]
\textbf{\large!} &
\textbf{Always justify in context.}\enspace
The AP exam expects a sentence explaining \emph{why} two events cannot
(or can) both occur, not just the label. \\[6pt]
\textbf{\large!} &
\textbf{Mutually exclusive $\neq$ Independent.}\enspace
These are different concepts. Mutually exclusive events with positive probabilities
are actually \emph{dependent} --- coming up in Lesson 4.6! \\
\end{tabularx}
\end{remindbox}

\end{document}
